% !TeX root = ../navier-stokes-manuscript.tex
% ============================================================
% 1.3 Gold Standard Proof
% ============================================================
\subsection{Gold Standard Proof}\label{sub:GoldStandardProof}

For smooth initial data, the standard whole-space local theory gives a smooth solution for a short time \parencite{Kato1984}, and the periodic problem has the corresponding local theory.
Let \([0,T_*)\) be the longest interval on which that solution remains smooth.
If \(T_*<\infty\), the local blow-up alternative forces every Sobolev norm above the continuation threshold to become unbounded as \(t\uparrow T_*\).
Keeping any one of those norms finite would continue the same solution past \(T_*\).
The Clay requirement asks the same solution to remain smooth through every finite time.
So the direct burden is to get a bound from \(u_0\), \(\nu\), and the equations that keeps the continuation norm finite on every finite interval.
A bound obtained this way is called an \emph{a priori estimate}.
Every estimate has to come from the same initial fluid and the same equations, without putting the needed bound in as an assumption.

To see what these estimates have to buy, start with a spatial Sobolev norm.
It measures the \(L^2\) size of the velocity together with a finite number of its spatial derivatives.
In three dimensions, Sobolev embedding gives \parencite[Theorem~7.28 for \(\mathbb R^3\)]{Cerda2010}
\[
  H^s(\Omega)\hookrightarrow C^k(\Omega)
  \qquad\text{whenever}\qquad
  s>k+\frac32,
  \qquad
  \Omega\in\{\mathbb R^3,\mathbb T^3\},
\]
where the periodic statement is the corresponding Fourier-series embedding.
One \(H^s\) bound buys only finitely many continuous spatial derivatives.
Fix \(s_0>5/2\).
An \(H^{s_0}\) bound gives \(C^1\) control, including a bounded gradient, but it still does not give \(C^\infty\).
The complete spatial target is control at every finite Sobolev order:
\[
  H^\infty(\Omega)
  :=
  \bigcap_{m=0}^{\infty}H^m(\Omega)
  \subset
  C^\infty(\Omega).
\]
The pressure and time derivatives still have to be recovered from the equations once that spatial control has been secured.
A natural first calculation starts with kinetic energy and asks whether its \(L^2\) control can be raised to one spatial derivative.
Work first on \(\mathbb T^3\), where no boundary terms appear; the same identity holds on \(\mathbb R^3\) by the usual cutoff-and-density argument for finite-energy smooth solutions.
Pair the momentum equation with \(u\) and integrate over the domain.
Incompressibility makes transport and pressure vanish from this particular total balance, while viscosity removes energy through spatial variation:
\[
  \begin{aligned}
  \int_{\mathbb T^3}(u\cdot\nabla)u\cdot u\,dx
  &=
  \frac12\int_{\mathbb T^3}u\cdot\nabla|u|^2\,dx
  =0,
  \\
  \int_{\mathbb T^3}\nabla p\cdot u\,dx
  &=
  -\int_{\mathbb T^3}p\,\nabla\cdot u\,dx
  =0,
  \\
  \nu\int_{\mathbb T^3}\Delta u\cdot u\,dx
  &=
  -\nu\lVert\nabla u\rVert_{L^2}^2.
  \end{aligned}
\]
The result is the exact energy identity
\[
  \frac12\frac{d}{dt}\lVert u(t)\rVert_{L^2}^2
  +
  \nu\lVert\nabla u(t)\rVert_{L^2}^2
  =0.
\]
After integration from \(0\) to \(t\), it becomes
\[
  \lVert u(t)\rVert_{L^2}^2
  +
  2\nu\int_0^t\lVert\nabla u(s)\rVert_{L^2}^2\,ds
  =
  \lVert u_0\rVert_{L^2}^2.
\]

The \(L^2\) size of the velocity stays bounded at every time, and viscosity pays for the total gradient cost accumulated over the interval.
Let
\[
  Y(t):=\lVert\nabla u(t)\rVert_{L^2}^2.
\]
The energy identity gives
\[
  2\nu\int_0^tY(s)\,ds
  \leq
  \lVert u_0\rVert_{L^2}^2.
\]
It does not give a bound on \(Y(t)\) at each individual time.
To control the gradient itself, the calculation has to be raised by one spatial derivative.

Pair the momentum equation with \(-\Delta u\).
The pressure term vanishes again because \(-\Delta u\) is divergence-free, while transport now survives:
\[
  \frac12Y'(t)
  +
  \nu\lVert\Delta u\rVert_{L^2}^2
  =
  \int_{\mathbb T^3}(u\cdot\nabla)u\cdot\Delta u\,dx.
\]
After integration by parts, the transport term is a product of three first derivatives:
\[
  \int_{\mathbb T^3}(u\cdot\nabla)u\cdot\Delta u\,dx
  =
  -\sum_{i,j,k=1}^3
  \int_{\mathbb T^3}
  \partial_k u_j\,\partial_j u_i\,\partial_k u_i\,dx.
\]
Hölder's inequality, interpolation between \(L^2\) and \(L^6\), the Sobolev control of the \(L^6\) factor, and Young's inequality give
\[
  \begin{aligned}
  \left|
    \int_{\mathbb T^3}(u\cdot\nabla)u\cdot\Delta u\,dx
  \right|
  &\leq C\lVert\nabla u\rVert_{L^3}^3
  \\
  &\leq C\lVert\nabla u\rVert_{L^2}^{3/2}
           \lVert\Delta u\rVert_{L^2}^{3/2}
  \\
  &\leq \frac{\nu}{2}\lVert\Delta u\rVert_{L^2}^2
           +C\nu^{-3}Y(t)^3.
  \end{aligned}
\]
Young's inequality has split the transport cost into one part viscosity can absorb and one part that is cubic in \(Y\).
After that absorption, the estimate becomes
\[
  Y'(t)
  +
  \nu\lVert\Delta u\rVert_{L^2}^2
  \leq
  C\nu^{-3}Y(t)^3.
\]
Integrating in time gives
\[
  Y(t)
  +
  \nu\int_0^t\lVert\Delta u(s)\rVert_{L^2}^2\,ds
  \leq
  Y(0)
  +
  C\nu^{-3}\int_0^tY(s)^3\,ds.
\]

The first estimate paid for \(\int_0^tY(s)\,ds\).
The next one asks for \(\int_0^tY(s)^3\,ds\).
A nonnegative function can have finite total area even when the area under its cube is infinite; tall, narrow spikes are enough to do it.
This does not show that a Navier\textendash Stokes solution forms those spikes.
It shows that the basic energy estimate carries no information that rules them out.
This calculation has stalled at the first derivative; it has not yet reached a continuation-grade \(H^{s_0}\) bound.

The full Sobolev ladder does not require a new global estimate at every order.
Assume for contradiction that \(T_*<\infty\) and suppose this one norm stays bounded all the way to the candidate endpoint:
\[
  \sup_{0\leq t<T_*}\lVert u(t)\rVert_{H^{s_0}}<\infty.
\]
Sobolev embedding then keeps \(\lVert\nabla u\rVert_{L^\infty}\) bounded, and hence integrable, on \([0,T_*)\).
For every finite \(m\geq s_0\), the standard higher-order estimate has the form
\[
  \frac12\frac{d}{dt}\lVert u\rVert_{H^m}^2
  +
  \nu\lVert\nabla u\rVert_{H^m}^2
  \leq
  C_m\lVert\nabla u\rVert_{L^\infty}\lVert u\rVert_{H^m}^2.
\]
The smooth initial datum already belongs to every \(H^m\), so Grönwall's inequality propagates each fixed finite order on the same interval.
The constants may depend on \(m\); no single constant has to control every derivative order at once.
Sobolev embedding then gives every continuous spatial derivative of \(u\).
This repeated climb from one secured strong bound through every higher order is the Sobolev bootstrap.
The local blow-up alternative would then continue the same solution past \(T_*\), contradicting its maximality.
A Gold estimate that supplies this bound at every candidate finite endpoint would force \(T_*=\infty\).

Once that global continuation has been secured, pressure is recovered from the same velocity, not added as a separate regularity assumption.
Taking the divergence of the momentum equation gives
\[
  -\Delta p
  =
  \sum_{i,j=1}^3\partial_i\partial_j(u_i u_j).
\]
On \(\mathbb T^3\), this relation fixes \(p\) by requiring zero spatial mean.
On \(\mathbb R^3\), the corresponding normalization is
\[
  p=R_iR_j(u_i u_j),
\]
where \(R_i\) is the \(i\)-th Riesz transform.
Elliptic estimates recover \(p\) at every spatial Sobolev order already available for \(u\).
On every closed interval inside the continued solution, the propagated Sobolev regularity remains continuous in time, so the momentum equation gives the first time derivative,
\[
  \partial_tu
  =
  \nu\Delta u
  -(u\cdot\nabla)u
  -\nabla p.
\]
Differentiating the pressure relation and the momentum equation repeatedly recovers every higher time derivative of both fields.
In exact Sobolev terms, this induction gives
\[
  \partial_t^r u,\ \partial_t^r p
  \in
  C\bigl([0,T];H^m(\Omega)\bigr)
  \qquad
  \text{for every }r,m\in\mathbb N_0
  \text{ and every finite }T>0.
\]
This is how the Sobolev estimates reach the spacetime \(C^\infty\) velocity\textendash pressure pair required above.

A Gold proof has to derive one continuation-grade bound from the original datum and the original equations that stays finite through every finite time.
The Sobolev bootstrap, the pressure relation, and the momentum equation then carry that bound into full smoothness.
For viscosity to \emph{pay the whole bill}, its dissipation has to control enough nonlinear transport to supply the missing bound.
In the calculation above, the energy identity pays for \(\int_0^tY(s)\,ds\), while the first differentiated estimate asks for \(\int_0^tY(s)^3\,ds\).
Finding a direct estimate that closes this gap, or reaches the continuation bound another way, without assuming the strong quantity it needs is the missing Gold step.
